\subsection{Null functions and equality almost everywhere}\label{sec:ae}

The seminorm $\|\cdot\|_I$ is in general not a norm.
We set up here the operation that collapses its kernel.
This quotient is used essentially in Section~\ref{chap:operator}.

\subsubsection{Null functions and the quotient}

\begin{definition}\label{def:null-function}
  Let $(X,L,I)$ be an integration system on a set $X$. We put
  \[
    \mathcal{N}_I:=\{f\in L\mid I(|f|)=0\}
  \]
  and call the functions belonging to $\mathcal{N}_I$ the \textbf{null functions}
  of $(X,L,I)$.
\end{definition}

\begin{lemma}\label{lem:null-is-ideal}
  $\mathcal{N}_I$ is an ideal of $L$.
  That is, $\mathcal{N}_I$ is a linear subspace, and if $f\in\mathcal{N}_I$ and
  $f'\in L$ satisfies $|f'|\le|f|$ then $f'\in\mathcal{N}_I$.
\end{lemma}
\begin{proof}
  $\mathcal{N}_I$ is the kernel of the seminorm $\|\cdot\|_I$, hence a linear
  subspace (Lemma~\ref{lem:l1-seminorm}).
  Moreover, if $|f'|\le|f|$ then
  Lemma~\ref{lem:integration-preserve-order} gives $I(|f'|)\le I(|f|)=0$, and
  positivity of $I$ gives $I(|f'|)=0$.
\end{proof}

\begin{definition}\label{def:null-quotient}
  We write the quotient space as
  \[
    \widetilde{L}:=L/\mathcal{N}_I
  \]
  and denote the quotient map by $\pi_I:L\to\widetilde{L}$.
\end{definition}

\begin{lemma}\label{lem:quotient-riesz}
  $\widetilde{L}$ carries the structure of a Riesz space for the order defined by
  declaring $\pi_I(f)\ge0$ to mean ``there is an $f'\in\mathcal{N}_I$ with
  $f+f'\ge0$'', and $\pi_I$ is a Riesz homomorphism.
  That is,
  \[
    \pi_I(f\vee g)=\pi_I(f)\vee\pi_I(g),
    \qquad
    \pi_I(|f|)=|\pi_I(f)| .
  \]
  Furthermore
  \[
    \widetilde{I}:\widetilde{L}\to\mathbb{R};\quad \pi_I(f)\mapsto I(f)
  \]
  is a well-defined positive linear functional, and
  $\|\pi_I(f)\|:=\widetilde{I}(|\pi_I(f)|)=\|f\|_I$ is a norm.
\end{lemma}
\begin{proof}
  That the quotient of a Riesz space by an ideal is a Riesz space and that the
  quotient map is a Riesz homomorphism is general theory (see for instance
  \cite{LuxemburgZaanen}).
  By Lemma~\ref{lem:null-is-ideal} this applies here.

  $\widetilde{I}$ is well defined because $f\in\mathcal{N}_I$ implies
  $|I(f)|\le I(|f|)=0$
  (Lemma~\ref{lem:triangle-ineq-for-integration}).
  Positivity and linearity follow from those of $I$.
  Finally
  $\widetilde{I}(|\pi_I(f)|)=\widetilde{I}(\pi_I(|f|))=I(|f|)=\|f\|_I$, and this
  vanishes only when $f\in\mathcal{N}_I$, that is, when $\pi_I(f)=0$.
\end{proof}

\begin{lemma}\label{lem:quotient-is-banach}
  If $(X,L,I)$ is a Daniell closed integration system then
  $(\widetilde{L},\|\cdot\|)$ is a Banach space.
\end{lemma}
\begin{proof}
  If $\{\pi_I(f_k)\}$ is Cauchy then $\{f_k\}$ is a Cauchy sequence with respect
  to $\|\cdot\|_I$.
  By Theorem~\ref{thm:l1-norm-complete} its limit $f\in L$ exists, and
  $\|\pi_I(f)-\pi_I(f_k)\|=\|f-f_k\|_I\to0$.
\end{proof}

\begin{lemma}[compatibility of the Daniell extension with the quotient]\label{lem:quotient-embedding}
  Let $(X,L,I)$ be an integration system on a set $X$ and let
  $(X,L^1,I^1)$ be its Daniell extension.
  Then
  \[
    \mathcal{N}_{I^1}\cap L=\mathcal{N}_I .
  \]
  Hence the inclusion $L\subset L^1$ induces an injective Riesz homomorphism
  \[
    j_I:\widetilde{L}\longrightarrow\widetilde{L^1}
  \]
  satisfying
  \[
    \pi_{I^1}|_L=j_I\circ\pi_I .
  \]
  Furthermore $j_I$ preserves the norm.
\end{lemma}
\begin{proof}
  By Lemma~\ref{lem:daniell-ext-contains-source} we have $|f|\in L$ and
  $I^1(|f|)=I(|f|)$ for every $f\in L$.
  Hence $I^1(|f|)=0$ and $I(|f|)=0$ are equivalent, and we obtain
  $\mathcal{N}_{I^1}\cap L=\mathcal{N}_I$.

  It follows that the kernel of $\pi_{I^1}|_L:L\to\widetilde{L^1}$ is
  $\mathcal{N}_I$, so an injective linear map
  $j_I:\widetilde{L}\to\widetilde{L^1}$ is induced, satisfying
  $\pi_{I^1}|_L=j_I\circ\pi_I$.
  Both $\pi_I$ and $\pi_{I^1}$ are Riesz homomorphisms
  (Lemma~\ref{lem:quotient-riesz}) and $\pi_I$ is surjective, so $j_I$ is a
  Riesz homomorphism as well.
  Finally, for $f\in L$,
  \[
    \|j_I(\pi_I(f))\|=\|\pi_{I^1}(f)\|=I^1(|f|)=I(|f|)=\|\pi_I(f)\| .
  \]
\end{proof}

\subsubsection{ae-equivalence of morphisms}

\begin{definition}\label{def:ae-equal}
  Let $(X,L,I)$ and $(Y,M,J)$ be integration systems on sets $X$ and
  $Y$ respectively.
  We say that two morphisms of integration systems
  \[
    F,G:(X,L,I)\to(Y,M,J)
  \]
  are \textbf{equal almost everywhere} if
  \[
    J\bigl(|F^*f-G^*f|\bigr)=0
  \]
  holds for every $f\in L$, that is, if $F^*f-G^*f\in\mathcal{N}_J$ holds for
  every $f\in L$.
  In this case we write
  \[
    F=_{\mathrm{ae}}G .
  \]
\end{definition}

\begin{lemma}\label{lem:ae-is-congruence}
  For the morphisms of integration systems between each pair of objects, the
  relation $=_{\mathrm{ae}}$ is an equivalence relation and is compatible with
  composition of morphisms.
  That is, if morphisms of integration systems
  \[
    F,F':(X,L,I)\to(Y,M,J),
    \qquad
    G,G':(Y,M,J)\to(Z,N,K)
  \]
  satisfy
  \[
    F=_{\mathrm{ae}}F',
    \qquad
    G=_{\mathrm{ae}}G' ,
  \]
  then
  \[
    G\circ F=_{\mathrm{ae}}G'\circ F' .
  \]
\end{lemma}
\begin{proof}
  Reflexivity and symmetry are immediate from the definition. As for
  transitivity, suppose morphisms of integration systems
  $F,F',F'':(X,L,I)\to(Y,M,J)$ satisfy $F=_{\mathrm{ae}}F'$ and
  $F'=_{\mathrm{ae}}F''$. For every $f\in L$,
  \[
    \begin{aligned}
      0
      &\le J\bigl(|F^*f-(F'')^*f|\bigr)\\
      &\le J\bigl(|F^*f-(F')^*f|\bigr)
        +J\bigl(|(F')^*f-(F'')^*f|\bigr)
      =0 ,
    \end{aligned}
  \]
  so $F=_{\mathrm{ae}}F''$.

  Next we show compatibility with composition. For every $f\in L$, since
  morphisms of Riesz systems preserve absolute values and morphisms of
  integration systems preserve the integral,
  \[
    \begin{aligned}
      &K\bigl(|G^*F^*f-(G')^*(F')^*f|\bigr)\\
      &\quad\le
      K\bigl(|G^*F^*f-G^*(F')^*f|\bigr)
      +K\bigl(|G^*(F')^*f-(G')^*(F')^*f|\bigr)\\
      &\quad=
      J\bigl(|F^*f-(F')^*f|\bigr)
      +K\bigl(|G^*(F')^*f-(G')^*(F')^*f|\bigr)
      =0.
    \end{aligned}
  \]
  Therefore $G\circ F=_{\mathrm{ae}}G'\circ F'$.
\end{proof}

\begin{definition}\label{def:int-ae}
  The category having the same objects as the category $\mathbf{Int}$ of
  integration systems and whose morphisms are the equivalence classes under
  $=_{\mathrm{ae}}$ is written
  \[
    \mathbf{Int}_{\mathrm{ae}}
  \]
  and called the \textbf{ae-category of integration systems}.
  By Lemma~\ref{lem:ae-is-congruence}, composition of morphisms is defined
  independently of the choice of representatives.

  Among the objects of $\mathbf{Int}_{\mathrm{ae}}$, we write the full
  subcategory consisting of the integration systems satisfying the monotone
  convergence theorem as
  \[
    \mathbf{MInt}_{\mathrm{ae}},
  \]
  and the full subcategory consisting of the Daniell closed integration systems
  as
  \[
    \mathbf{DInt}_{\mathrm{ae}} .
  \]
\end{definition}
